\documentclass[11pt]{article}
\usepackage[letterpaper,margin=1in]{geometry}
\usepackage{amsmath,amssymb,amsthm,booktabs,microtype,xurl}
\usepackage[hidelinks]{hyperref}
\newtheorem{theorem}{Theorem}
\newtheorem{lemma}[theorem]{Lemma}
\newtheorem{corollary}[theorem]{Corollary}
\newcommand{\C}{\mathbb C}
\newcommand{\Q}{\mathbb Q}
\newcommand{\SIC}{\operatorname{SIC}}

\title{A 14-variable polynomial map with everywhere\\
unipotent Jacobian and a three-point fiber}
\author{Alec Kriebel\\[2pt]\large with heavy assistance from ChatGPT 5.6 Sol}
\date{Branch draft prepared 22 July 2026\\Research draft; not peer reviewed}

\begin{document}
\maketitle

\begin{abstract}
We give an exact rational polynomial vector $g\in\Q[Z]^{14}$, with
24 monomials, such that
\[
 \det(I+sJg)=1
\]
identically, and such that $I+g$ has three displayed rational points in one
fiber.  Generically, $Jg$ is regular nilpotent, with Jordan type $(14)$.
The generic degree is three and the geometric monodromy is $S_3$.
For the image operator
$E(\xi^\alpha p(Z))=\partial_Z^\alpha p(Z)$, put
$A=-\sum_{j=1}^{14}\xi_jg_j$ and $b=x+y+u_{11}$.  We prove
\[
 E(A^m)=0,\qquad E(bA^m)\ne0\qquad(m\geq1).
\]
Thus $\SIC(14)$ fails, with an obstruction at every exponent.  We also
record degree-seven homogeneous and Hessian-nilpotent companions.  The
dimension and term counts are proved optimal only in a precisely stated
constant-state realization ansatz.
\end{abstract}

\noindent\textbf{Verification disclaimer.}
The human author is a complete amateur exploring the limits of AI-assisted
mathematics and cannot independently verify these claims.  This is a research
artifact for expert checking, not an established result.  Priority and
minimality have only the narrow scope stated below.

\section{The three-variable source}

Put $u=1+xy$ and
\[
\begin{aligned}
F_1&=u^3z+y^2u(4+3xy),\\
F_2&=y+3xu^2z+3xy^2(4+3xy),\\
F_3&=2x-3x^2y-x^3z.
\end{aligned}
\]
Use the identity-linear normalization
\begin{equation}
 \Phi=(F_3/2,F_2,F_1)=X+H_2+H_3+\cdots+H_7. \label{eq:Phi}
\end{equation}
The nonzero homogeneous pieces are displayed in Table~\ref{tab:pieces}.
\begin{table}[ht]
\centering
\begin{tabular}{c|ccc}
\toprule
$d$&$H_{d,1}$&$H_{d,2}$&$H_{d,3}$\\
\midrule
2&0&$3xz$&$4y^2$\\
3&$-\tfrac32x^2y$&$12xy^2$&$3xyz$\\
4&$-\tfrac12x^3z$&$6x^2yz$&$7xy^3$\\
5&0&$9x^2y^3$&$3x^2y^2z$\\
6&0&$3x^3y^2z$&$3x^2y^4$\\
7&0&0&$x^3y^3z$\\
\bottomrule
\end{tabular}
\caption{Homogeneous pieces of $\Phi-I$.}
\label{tab:pieces}
\end{table}

Direct differentiation gives $\det J\Phi=1$.  The three points
\begin{equation}
 r_0=(0,0,-\tfrac14),\quad
 r_+=(1,-\tfrac32,\tfrac{13}2),\quad
 r_-=(-1,\tfrac32,\tfrac{13}2) \label{eq:sourcepoints}
\end{equation}
all map to $(0,0,-1/4)$.

\section{The 14-variable construction}

Take state coordinates
\[
 U=(u_{11},u_{12};u_{21},u_{22},u_{23},u_{24};
 u_{31},u_{32},u_{33},u_{34},u_{35}).
\]
Let $B\in\operatorname{Mat}_{3\times11}(\Q)$ select the three chain heads,
so $BU=(u_{11},u_{21},u_{31})^T$.  Let
$N=J_2\oplus J_4\oplus J_5$, with ones immediately above the diagonal in
each Jordan block, and put
\[
C=\begin{pmatrix}
H_{3,1}\\H_{4,1}\\
H_{3,2}\\H_{4,2}\\H_{5,2}\\H_{6,2}\\
H_{3,3}\\H_{4,3}\\H_{5,3}\\H_{6,3}\\H_{7,3}
\end{pmatrix}.
\]
Define
\begin{equation}
 g(X,U)=\bigl(H_2(X)+BU,-C(X)-NU\bigr). \label{eq:blockg}
\end{equation}
In the coordinate order $Z=(x,y,z,U)$, its components are
\begingroup\small
\begin{align}
g_1&=u_{11},&
g_2&=3xz+u_{21},&
g_3&=4y^2+u_{31},\nonumber\\
g_4&=\tfrac32x^2y-u_{12},&
g_5&=\tfrac12x^3z,\nonumber\\
g_6&=-12xy^2-u_{22},&
g_7&=-6x^2yz-u_{23},\nonumber\\
g_8&=-9x^2y^3-u_{24},&
g_9&=-3x^3y^2z,\nonumber\\
g_{10}&=-3xyz-u_{32},&
g_{11}&=-7xy^3-u_{33},\nonumber\\
g_{12}&=-3x^2y^2z-u_{34},&
g_{13}&=-3x^2y^4-u_{35},\nonumber\\
g_{14}&=-x^3y^3z. \label{eq:gdisplay}
\end{align}
\endgroup
There are exactly 24 monomials.  Hence
\begin{equation}
 A(\xi,Z)=-\sum_{j=1}^{14}\xi_jg_j(Z) \label{eq:A}
\end{equation}
has 24 monomials in 28 indeterminates and total degree eight.

\section{Determinant pencil and regular nilpotency}

\begin{theorem}\label{thm:main}
For the vector in~\eqref{eq:gdisplay},
\[
 \det(I+sJg)=1.
\]
The generic Jordan type of $Jg$ is $(14)$.  The map $I+g$ has three
distinct rational points in one fiber, has generic degree three, and has
geometric monodromy $S_3$.
\end{theorem}

\begin{proof}
The Jacobian pencil has block form
\[
 I+sJg=
 \begin{pmatrix}
 I_3+sJH_2&sB\\
 -sJC&I_{11}-sN
 \end{pmatrix}.
\]
Since $N^5=0$,
\[
 \det(I-sN)=1,\qquad
 (I-sN)^{-1}=\sum_{k=0}^4s^kN^k.
\]
The chain definition gives $BN^kC=H_{k+3}$ for $0\leq k\leq4$, with
absent components zero.  The Schur complement is therefore
\[
\begin{aligned}
I_3+sJH_2+s^2B(I-sN)^{-1}JC
 &=I_3+\sum_{d=2}^7s^{d-1}JH_d(X)\\
 &=J\Phi(sX).
\end{aligned}
\]
Its determinant is one, proving the pencil identity.  Cayley--Hamilton gives
$(Jg)^{14}=0$, while exact multiplication gives
\begin{equation}
 \bigl((Jg)^{13}\bigr)_{1,5}=-3x^6y^4z\ne0. \label{eq:index}
\end{equation}
The generic nilpotency index is therefore 14, which forces the single Jordan
block $(14)$.

For a source point $r$, put $U(r)=(I-N)^{-1}C(r)$.  Then
\[
 (I+g)(r,U(r))=(\Phi(r),0).
\]
Thus the three points in~\eqref{eq:sourcepoints} lift to the explicit points
displayed in~\eqref{eq:liftedpoints} below and have common image.

Finally, for a generic target $(Y,V)$ the state equations give
\[
 U=(I-N)^{-1}(V+C(X)),\qquad
 \Phi(X)=Y-B(I-N)^{-1}V.
\]
Thus $I+g$ is stably equivalent to $\Phi\times I_{11}$.  The exact cubic
resolvent, irreducibility, and nonsquare-discriminant certificates for
$\Phi$ are given in Discovery 04; they show generic degree three and
geometric monodromy $S_3$.
\end{proof}

The lifted points promised in the proof are
\begingroup\small
\begin{align}
p_0={}&(0,0,-\tfrac14;0,0;0,0,0,0;0,0,0,0,0),\nonumber\\
p_+={}&(1,-\tfrac32,\tfrac{13}2;-1,-\tfrac{13}4;
-18,-45,\tfrac{27}2,\tfrac{351}8;\nonumber\\
&\hspace{26mm}-\tfrac{63}4,\tfrac{27}2,\tfrac{297}8,
-\tfrac{27}4,-\tfrac{351}{16}),\nonumber\\
p_-={}&(-1,\tfrac32,\tfrac{13}2;1,\tfrac{13}4;
18,45,-\tfrac{27}2,-\tfrac{351}8;\nonumber\\
&\hspace{26mm}-\tfrac{63}4,\tfrac{27}2,\tfrac{297}8,
-\tfrac{27}4,-\tfrac{351}{16}). \label{eq:liftedpoints}
\end{align}
\endgroup
Their common image is $(0,0,-1/4,0,\ldots,0)$.

For every nonzero scalar $s$, the same elimination applied to $I+sg$ gives
\[
 s^{-1}\Phi(sX)=Y-sB(I-sN)^{-1}V.
\]
Consequently every nonzero pencil member is noninjective and retains generic
degree three and monodromy $S_3$.

\section{Consequences for injectivity conjectures}

The derivative of $I+g$ is everywhere unipotent, but
\eqref{eq:liftedpoints} proves noninjectivity.  This gives an explicit
dimension-14 counterexample to Campbell's Unipotent Jacobian Univalence
Conjecture and to the all-eigenvalues-one form of Chamberland's injectivity
conjecture.

It also gives a concrete failure of Kulikov's $JN(14)$.  For
$c\in\C^{14}$ set
\[
 N_c(Z)=c-g(Z).
\]
Then $JN_c=-Jg$ is nilpotent, and fixed points of $N_c$ are exactly the
points of $(I+g)^{-1}(c)$.  The displayed value gives three rational fixed
points; a generic value gives three geometric fixed points.

These are algebraic injectivity statements.  They should not be confused
with the classical dynamical Markus--Yamabe conjecture, which has older
counterexamples, or with valid theorems using singular-value rather than
eigenvalue bounds.

\section{The Special Image Conjecture}

For $n$ pairs of variables define
\[
 E_n:\C[\xi_1,\ldots,\xi_n,Z_1,\ldots,Z_n]\longrightarrow\C[Z],
 \qquad E_n(\xi^\alpha p)=\partial_Z^\alpha p.
\]
The Special Image Conjecture says that $\ker E_n$ is a Mathieu--Zhao
subspace.

\begin{lemma}[scalar-parameter inversion]\label{lem:inverse}
Suppose $\det(I+tJg)=1$, put $A=-\xi\cdot g$, and let $Q_t$ be the formal
inverse of $I+tg$.  For every polynomial $p$,
\begin{equation}
 p(Q_t(Z))=\sum_{m\geq0}\frac{t^m}{m!}E_n(pA^m). \label{eq:AG}
\end{equation}
\end{lemma}

\begin{proof}
Apply the Abhyankar--Gurjar inversion formula to
$I+tg=I-H_t$, where $H_t=-tg$.  The determinant factor is one, and the
multinomial theorem collects terms of total $\xi$-degree $m$.  The formula
is valid $t$-adically: each coefficient receives only finitely many
contributions.
\end{proof}

\begin{theorem}\label{thm:SIC}
For $A$ in~\eqref{eq:A} and $b=x+y+u_{11}$,
\[
 E_{14}(A^m)=0,\qquad E_{14}(bA^m)\ne0
 \qquad(m\geq1).
\]
Consequently $\SIC(14)$ is false.
\end{theorem}

\begin{proof}
Taking $p=1$ in~\eqref{eq:AG} gives $E_{14}(A^m)=0$ for every $m\geq1$.
For the sharp obstruction, specialize the inverse target to
\[
 Z_*=(\tfrac12,0,1,0,\ldots,0).
\]
Let $\tau\in t\Q[[t]]$ be the unique solution of
\begin{equation}
 \tau+t\tau^3=\frac t2. \label{eq:tau}
\end{equation}
State elimination followed by the source cubic reconstruction gives
\[
 Q_{t,x}(Z_*)=\tau',\qquad
 Q_{t,y}(Z_*)=-3\tau,\qquad
 Q_{t,u_{11}}(Z_*)=-\frac{\tau'-1/2}{t}.
\]
The last identity is also immediate from the first target equation
$Q_x+tQ_{u_{11}}=1/2$.

Write $b(Q_t)(Z_*)=\sum_{m\geq0}q_mt^m$.  Lagrange inversion
in~\eqref{eq:tau} gives
\begin{align}
q_{3k}&=(-1)^k\frac{\binom{3k+1}{k}}{2^{2k+1}},\nonumber\\
q_{3k+1}&=(-1)^{k+1}
 \frac{3\binom{3k+1}{k}}{(3k+1)2^{2k+1}},\nonumber\\
q_{3k+2}&=(-1)^k\frac{\binom{3k+4}{k+1}}{2^{2k+3}}. \label{eq:q}
\end{align}
Every number in~\eqref{eq:q} is nonzero.  Comparing coefficients
in~\eqref{eq:AG} yields
\[
 E_{14}(bA^m)(Z_*)=m!q_m\ne0
\]
for every $m\geq1$.
\end{proof}

\section{Optimality inside the state ansatz}

No global minimality claim is made.  Consider only realizations
\begin{equation}
 \widetilde g(X,U)=(H_2(X)+\widetilde BU,
 -\widetilde C(X)-\widetilde NU) \label{eq:ansatz}
\end{equation}
with constant matrices, nilpotent $\widetilde N$, and exact Markov parameters
\[
 \widetilde B\widetilde N^k\widetilde C=H_{k+3}\qquad(k\geq0).
\]
Order the eleven distinct tail monomials as
\[
x^2y,\ xy^2,\ xyz,\ x^3z,\ x^2yz,\ xy^3,\ x^2y^3,\ x^2y^2z,\
x^3y^2z,\ x^2y^4,\ x^3y^3z.
\]
Let $M_k\in\operatorname{Mat}_{3\times11}(\Q)$ be the coefficient matrix
of $H_{k+3}$ in these columns, with $M_k=0$ for $k\geq5$.  For
$0\leq j\leq5$, form
$\mathcal H_j=(M_{j+r+c})_{0\leq r,c<5-j}$.  Their exact ranks are
\[
 11,8,5,3,1,0.
\]
For $\mathcal H_0$, in standard block order, rows
$1,2,3,4,5,6,8,9,11,12,15$ and columns $1,\ldots,11$ give the exact minor
\[
 275562=2\cdot3^9\cdot7\ne0.
\]
Every realization factors $\mathcal H_0$ through its state space, so at
least eleven states are required.  The shifted matrix $\mathcal H_1$ factors
through $\widetilde N$, and hence its rank gives
$\operatorname{rank}\widetilde N\geq8$.  Thus dimension 14 is minimal
in~\eqref{eq:ansatz}.  Moreover
$\operatorname{rank}\widetilde B\geq3$,
and $\widetilde C$ must contain all eleven monomial inputs.  Together with
the two terms of $H_2$, this gives
\[
 2+3+8+11=24
\]
as a term lower bound.  The displayed construction attains it.

\section{Homogeneous companions}

Introduce $w$ and homogenize to degree seven:
\begin{equation}
 h(Z,w)=\bigl(w^7g(Z/w),0\bigr)\in\Q[Z,w]^{15}. \label{eq:h}
\end{equation}
This vector has 24 terms, and its zero last row gives
\[
 \det(I+sJh)=\det(I+sw^6Jg(Z/w))=1.
\]
Thus $Jh$ is nilpotent.  The points $(p_0,1),(p_+,1),(p_-,1)$ collide
under $I+h$, while degree and monodromy remain $3$ and $S_3$.  The exact
checker certifies $(Jh)^{13}\ne0$, so its generic nilpotency index is at
least 14.  We make no exact index assertion without a separate degree-14
zero certificate.

As a standard consequence of the de Bondt--van den Essen symmetrization
lemma, the polynomial
\begin{equation}
 P(A,B)=i\sum_{j=1}^{15}h_j(A+iB)B_j \label{eq:P}
\end{equation}
is a homogeneous degree-eight Hessian-nilpotent polynomial in 30 variables.
The usual linear solve with $I+Jh(x)^T$ transports two displayed colliding
points to a collision of $I-\nabla P$.  This paragraph invokes the standard
lemma; the independently checked headline artifact is the 14-variable map.

\section{Scope and priority}

The state-space count
\[
 3+(4-2)+(6-2)+(7-2)=14
\]
is classical degree bookkeeping in Kulikov's reduction.  The abstract
existence of a 14-dimensional reduction is not claimed as new.  The asserted
contribution is the sparse explicit realization, its regular-nilpotent and
three-point certificates, and the consequence cascade, particularly the
every-exponent $\SIC(14)$ formula.

At the audit cutoff, no earlier public source was located giving this exact
object or an explicit $\SIC(n)$ witness for $n\leq14$.  That is provisional
evidence, not proof of worldwide priority.  No global smallest-dimension or
smallest-term claim is made.

\begin{thebibliography}{9}
\bibitem{Campbell}
L. A. Campbell,
\emph{Unipotent Jacobian Matrices and Univalent Maps},
\url{https://arxiv.org/abs/math/9907157}.

\bibitem{Kulikov}
V. S. Kulikov,
\emph{Jacobian Conjecture and Nilpotent Mappings},
\url{https://arxiv.org/abs/math/9803143}.

\bibitem{Zhao}
W. Zhao,
\emph{Images of Commuting Differential Operators of Order One with
Constant Leading Coefficients},
J. Algebra 324 (2010), 231--247;
\url{https://arxiv.org/abs/0902.0210}.

\bibitem{BondtEssen}
M. de Bondt and A. van den Essen,
\emph{A reduction of the Jacobian conjecture to the symmetric case},
Proc. Amer. Math. Soc. 133 (2005), 2201--2205.

\bibitem{Discovery04}
A. Kriebel, with ChatGPT 5.6 Sol,
\emph{Full wreath-product monodromy through the third iterate of an explicit
Keller map}, Discovery 04 (2026).
\end{thebibliography}

\end{document}
